\section{A sharp three-family theorem for compressed surface covers}
\label{sec:dihedral-covers}

Normal coordinates can encode far more product pieces than can be written
down individually.  A useful algorithmic interface therefore returns
families with binary multiplicities and answers queries about individual
members on demand.  We give an exact instance of this principle for
surface covers with dihedral monodromy.  The result is independent of the
chain-complex transfer algorithm: it is a geometric kernel for a
restricted class of possible parallelity-bundle presentations.

The underlying covering-space and voltage-graph principles are classical.
Permutation voltage assignments were developed by Gross and Tucker
\cite{GrossTucker1977}; their later paper already gives algorithms for
component counting and orientability \cite{GrossTucker1979}.  The present
contribution is an explicit binary-arithmetic specialization, including a
sharp bound of three component families, complete boundary-lift profiles,
and an independently checked implementation.  We do not claim a new
general covering-space classification or a replacement for the
Agol--Hass--Thurston algorithm \cite{AHT2006}.

\subsection{The input promise}

Let $F$ be a compact connected surface with $b\geq 1$ boundary components.
Fix a basepoint and free generators $z_1,\ldots,z_r$ for its fundamental
group, where $r=1-\chi(F)$.  Specify a $W$-sheeted covering by its
monodromy on the fibre $\mathbb Z/W\mathbb Z$:
\[
 g_i(x)=\epsilon_i x+a_i\pmod W,
 \qquad \epsilon_i\in\{1,-1\},\quad 0\leq a_i<W.
\]
The integer $W$ and the shifts $a_i$ are encoded in binary.  Because the
fundamental group is free, every such assignment defines a covering of
$F$.  It need not be connected.  Let $w_i\in\mathbb Z/2\mathbb Z$ be the
orientation character of $z_i$, and let $\beta_1,\ldots,\beta_b$ be based
words representing the boundary circles.  Their evaluated monodromies
are denoted by
\[
 h_j(x)=s_jx+c_j\pmod W.
\]
These words must be an actual peripheral system for the specified
surface.  Arbitrary words with zero orientation character are not by
themselves a certificate of that fact.

The delivered interface makes this promise concrete using a canonical
surface schema.  An orientable surface of genus $g$ has free generators
\[
 a_1,b_1,\ldots,a_g,b_g,c_1,\ldots,c_{b-1},
\]
all with orientation character zero.  The first $b-1$ boundary words are
$c_1,\ldots,c_{b-1}$, and the final one is the inverse of
\[
 [a_1,b_1]\cdots[a_g,b_g]c_1\cdots c_{b-1}.
\]
For a nonorientable surface with $g$ crosscaps, the generators are
$x_1,\ldots,x_g,c_1,\ldots,c_{b-1}$; the $x_i$ have orientation
character one and the $c_j$ have character zero.  The last boundary word
is the inverse of $x_1^2\cdots x_g^2c_1\cdots c_{b-1}$.
The disk is included, with no generators and one empty boundary word.
Closed surfaces are outside this implementation's interface.

\begin{theorem}[Three component families]
\label{thm:dihedral-three}
For the covering described above, all connected components can be
represented by at most three records.  Each record specifies a binary
multiplicity and one isomorphism class of connected covering of $F$.
The normalized input monodromy is stored once, shared by the records;
each record additionally supplies a representative residue orbit.
In particular, it determines the covering degree, orientability, genus
(or crosscap number), and the complete degree profile of the lifted
circles over each boundary component of $F$.

There is an algorithm constructing these records in time polynomial in
the input size, including $\log(W+1)$.  It also answers, without sheet
expansion, which covering component contains a specified fibre point
and which lifted boundary circle contains that point with respect to a
specified boundary word.  The bound of three is sharp, even when $F$ is
orientable.
\end{theorem}

We establish the theorem through explicit formulas.  The formulas also
specify the implementation completely.

\subsection{Translation orbits and the two exceptional residues}

If there are no generators of sign $-1$, set
\[
 d=\gcd(W,a_1,\ldots,a_r),\qquad m=W/d.
\]
If a negative-sign generator exists, fix one, written
\[
 R(x)=-x+a_0.
\]
Form a list of translation increments consisting of all $a_i$ with
$\epsilon_i=1$ and all $a_i-a_0$ with $\epsilon_i=-1$.  Reduce these
increments modulo $W$ and call them $u_1,\ldots,u_t$.  Set
\begin{equation}
 d=\gcd(W,u_1,\ldots,u_t),\qquad m=W/d.
 \label{eq:dihedral-gcd}
\end{equation}
An empty list has $d=W$.  The translation subgroup generated by the
monodromy consists exactly of translations by multiples of $d$.
Indeed, the product of two reflections has increment equal to the
difference of their shifts, and conjugating a translation by a
reflection reverses its increment.  Thus the stated increments generate
every translation, and Bezout's identity gives the converse inclusion.

Without a reflection, the connected components correspond to the $d$
residue classes modulo $d$, each containing $m$ sheets.  With a
reflection, pass to the quotient by translations.  The remaining action
is the involution
\[
 \iota(r)=a_0-r\pmod d.
\]
Its fixed residues satisfy
\begin{equation}
 2r\equiv a_0\pmod d.
 \label{eq:dihedral-fixed}
\end{equation}
Let $f$ be the number of such residues.  Explicitly, with
$e=\gcd(2,d)$,
\[
 f=\begin{cases}e,&e\mid a_0,\\0,&e\nmid a_0.\end{cases}
\]
There are $(d-f)/2$ components comprising pairs of residue classes,
each of degree $2m$, and $f\leq 2$ exceptional components, each of
degree $m$.  In particular, the component count is $(d+f)/2$.

The generic components are isomorphic as covers of $F$, not merely as
abstract surfaces.  For two nonfixed residues $r,s$, define the map on
their two-residue orbits by
\begin{align*}
 r+jd&\longmapsto s+jd,\\
 a_0-r+jd&\longmapsto a_0-s+jd
 \qquad (j\in\mathbb Z/m\mathbb Z).
\end{align*}
The two source residue classes are disjoint, so this is well-defined.
It commutes with translations by $d$ and with $R$, hence with every
generator of monodromy.  An equivariant bijection of fibres induces an
isomorphism of the corresponding connected covers
\cite[Section~1.3]{Hatcher2002}.  In the purely
translational case, adding $s-r$ gives the analogous equivariant
bijection between any two components.  There is therefore one generic
record and at most two exceptional records.  This argument remains
valid for $W=1,2$, even though distinct formal affine signs can then
describe the same permutation.

\begin{corollary}[Exact number of cover-isomorphism types]
\label{cor:dihedral-exact-types}
If no reflection is present, there is exactly one isomorphism type of
connected covering component.  If a reflection is present, the exact
number is
\[
 N_{\mathrm{iso}}=\mathbf 1_{\{d>f\}}+
 \begin{cases}
 0,&f=0,\\
 1,&f=1\text{ or }(f=2\text{ and }m\text{ is odd}),\\
 2,&f=2\text{ and }m\text{ is even}.
 \end{cases}
\]
In particular, when there are two fixed residues, their connected
$m$-sheeted covers are isomorphic over $F$ exactly when $m$ is odd.
\end{corollary}

\begin{proof}
Only the relation between two exceptional components remains to be
decided.  Such components exist precisely when $d$ and $a_0$ are even.
Their residues can be written $r$ and $s=r+d/2$, with representatives
chosen in $\{0,\ldots,d-1\}$.

Any covering isomorphism induces a fibre bijection commuting with
monodromy, hence with translation by $d$.  That translation is transitive
on each exceptional residue class.  The bijection is therefore forced
to have the form $x\mapsto x+A$ on the source class, where
\[
 A=s-r+dt=d/2+dt
 \qquad (t\in\mathbb Z/m\mathbb Z).
\]
Commuting also with the reference reflection requires
$2A=0\pmod W$.  Since $W=dm$, this is equivalent to
\[
 2t=-1\pmod m,
\]
which is solvable exactly when $m$ is odd.  Conversely, any such
translation commutes with every translation and reflection in the
monodromy group, so it supplies the required covering isomorphism.
For odd $m$, the central translation $x\mapsto x+W/2$ is one explicit
choice and swaps the two exceptional residues.

The generic components have degree $2m$, so cannot be isomorphic as
covers of $F$ to an exceptional component of degree $m$.  All generic
components are mutually isomorphic by the preceding proof.  These facts
give the formula.  It includes $W=1,2$: in the two-exception case with
$W=2$ one has $d=2,m=1$, and the two singleton fibres are exchanged by
translation by one.
\end{proof}

This distinction is visible even for two-sheeted connected covers.
Take a pair of pants, $W=4$, and generator permutations $x\mapsto x+2$
and $x\mapsto-x$.  Its two covering components are both four-holed
spheres.  On one fibre orbit the second generator is the identity,
whereas on the other it is a transposition.  The covers are therefore
nonisomorphic over the specified base, in agreement with $d=m=2$.

The orbit and cover-isomorphism conclusions use only the affine
monodromy action.  They therefore hold for coverings of any connected
CW complex with such monodromy.  The surface hypothesis is needed for
the orientation, peripheral, and genus interpretations below; a
presentation on a different base must still specify a genuine
monodromy action satisfying that base's relations.

\subsection{Orientability from a parity calculation}

The permutation action alone does not determine orientability when the
base surface is nonorientable.  One must retain the orientation bits of
the words producing the permutations.

Attach a bit $v_j$ to each translation increment $u_j$.  For a
positive-sign generator its bit is its orientation character.  For the
difference of the shifts of a reflection $g_i$ and the reference
reflection $R$, its bit is $w_i+w_0$ modulo two.  Include also the
relation $(W,0)$.  Bezout's algorithm produces integers with
\[
 d=k_0W+\sum_j k_ju_j.
\]
Define
\[
 \tau=\sum_j k_jv_j\pmod2.
\]
Check the following finite system:
\begin{equation}
 m\tau=0\pmod2,
 \qquad
 v_j=(u_j/d)\tau\pmod2\quad\hbox{for every }j.
 \label{eq:dihedral-parity}
\end{equation}

\begin{lemma}[Orientation criterion]
\label{lem:dihedral-orientation}
If \eqref{eq:dihedral-parity} fails, every covering component is
nonorientable.  If it holds, every component without a fixed residue is
orientable.  A component corresponding to a fixed residue $r$ is
orientable exactly when
\begin{equation}
 w_0+\frac{2r-a_0}{d}\tau=0\pmod2.
 \label{eq:dihedral-exception-orientation}
\end{equation}
In the absence of reflections, all components are orientable if and
only if \eqref{eq:dihedral-parity} holds.
\end{lemma}

\begin{proof}
Retain the word orientation in the lifted action on
$(\mathbb Z/W\mathbb Z)\times(\mathbb Z/2\mathbb Z)$.  A translation
acts by $(x,q)\mapsto(x+u,q+v)$, and the reference reflection acts by
$(x,q)\mapsto(-x+a_0,q+w_0)$.  A connected component is nonorientable
precisely when some word fixes a sheet and changes the orientation
coordinate.

The integer subgroup generated by $(W,0)$ and $(u_j,v_j)$ contains
$(d,\tau)$ by the displayed Bezout identity.  If any equality in
\eqref{eq:dihedral-parity} fails, subtracting the appropriate multiple
of $(d,\tau)$ from $(W,0)$ or $(u_j,v_j)$ gives $(0,1)$.
This is an orientation-reversing word fixing every sheet.  Conversely,
if all equalities hold, the signed translation subgroup is exactly
\[
 \{(jd,j\tau):j\in\mathbb Z/m\mathbb Z\}.
\]
The first equality ensures that the second coordinate is well-defined
modulo $m$.  Thus a translation fixing a sheet cannot reverse
orientation.

Every remaining signed group element is a translation followed by the
reference reflection.  Such an element fixes $x$ exactly when
\[
 jd=2x-a_0\pmod W.
\]
There is no solution on a nonfixed residue class.  On a fixed class,
write $x=r+yd$.  Then
\[
 j=\frac{2r-a_0}{d}+2y\pmod m.
\]
The orientation bit of this stabilizing reflection is consequently
$w_0+((2r-a_0)/d)\tau$.  It is independent of $y$, and independent of
the choice of $j$ modulo $m$ by \eqref{eq:dihedral-parity}.  This proves
all the assertions.
\end{proof}

It is insufficient to take a gcd of permutation shifts and ignore the
orientation character.  For example, over a M\"obius band, translation
by one on an even number of sheets gives a connected annulus cover;
the same construction with an odd number of sheets gives a connected
M\"obius-band cover.  The relation $m\tau=0$ distinguishes the cases.
For an orientable base, all orientation bits are zero and the criterion
automatically declares every component orientable, as it should.

\subsection{Boundary lift degrees and surface topology}

For one base boundary circle, evaluate its monodromy as
$h(x)=s x+c$.  A lifted boundary circle corresponds to a permutation
cycle of $h$ within a covering component; the cycle length is the degree
of that circle over the base boundary.  The full profile follows from
the following cases.

If $s=1$, then $d\mid c$.  Set
\[
 q=\gcd(m,c/d).
\]
On a component consisting of one residue class, there are $q$ boundary
lifts, each of degree $m/q$.  On a generic two-residue component there
are $2q$ such lifts, each of the same degree $m/q$.

If $s=-1$, a generic two-residue component contains $m$ cycles of
length two: the reflection exchanges its two residue classes and fixes
no sheet.  For a fixed residue $r$, write $x=r+dy$.  The reflection
acts on $y\in\mathbb Z/m\mathbb Z$ by
\[
 y\longmapsto\frac{c-2r}{d}-y.
\]
Thus, putting
\begin{equation}
 \nu=\begin{cases}
 \gcd(2,m),&\gcd(2,m)\mid(c-2r)/d,\\
 0,&\text{otherwise},
 \end{cases}
 \label{eq:dihedral-boundary-fixed}
\end{equation}
there are $\nu$ degree-one boundary lifts and $(m-\nu)/2$ degree-two
lifts.  Terms with multiplicity zero are omitted.  Every boundary
therefore needs at most two degree/multiplicity pairs per component
family.

Let $B_C$ be the resulting number of boundary circles of a component
$C$, and let $n_C$ be its covering degree.  Euler characteristic
multiplicativity gives
\[
 \chi(C)=n_C\chi(F).
\]
Its genus is determined without constructing a triangulation:
\[
 g(C)=\frac{2-B_C-n_C\chi(F)}2
 \quad\text{if }C\text{ is orientable},
\]
and its crosscap number is $2-B_C-n_C\chi(F)$ otherwise.  These
quantities are integral with the required signs because the input
specifies an actual surface covering.  The implementation checks these
identities as internal consistency conditions.

For the canonical interval bundle with orientable total space over
$C$, orientability of $C$ also determines the bundle type: it is the
product bundle when $C$ is orientable, and the twisted orientation
interval bundle otherwise.  A different independently prescribed
interval-bundle character would require running the same stabilizer
calculation with that character; the code does not infer such extra
data from a surface's topology.

\subsection{Sharpness, including orientable bases}

Take an orientable genus-one surface with two boundary circles, with
free generators $a,b,c$ and boundary words $c$ and $([a,b]c)^{-1}$.
Let $d\geq4$ and $m\geq2$ both be even, set $W=dm$, and assign
\[
 a(x)=x+d,\qquad b(x)=-x,\qquad c(x)=-x\pmod W.
\]
The translation divisor is $d$.  The fixed residues are $0$ and $d/2$.
The two boundary monodromies are reflections with shifts $0$ and
$2d$ up to the harmless sign determined by the path-composition
convention.  Applying the preceding formulas gives
\[
\begin{array}{c|c|c|c|c}
\text{family}&\text{multiplicity}&\text{degree}&\text{genus}&
 \text{boundary circles}\\\hline
\text{paired residues}&(d-2)/2&2m&m+1&2m\\
 r=0&1&m&m/2&m+2\\
 r=d/2&1&m&m/2+1&m
\end{array}
\]
All three rows occur and their genera are distinct.  They are therefore
three distinct homeomorphism types and, a fortiori, three distinct
cover-isomorphism types.  This proves sharpness in
Theorem~\ref{thm:dihedral-three}.  For example, $W=16,d=m=4$ yields
surfaces of genera $5,2,3$, with $8,6,4$ boundary circles, respectively.

\subsection{Marked queries and the output contract}

At the basepoint, a sheet $x$ belongs to the component indexed by
$r=x\bmod d$ if there is no reflection.  With a reflection, its
component is indexed by the unordered pair
\[
 \{r,a_0-r\bmod d\},
\]
with a singleton used for a fixed residue.  The function
\texttt{component\_key} returns precisely this label.

For a specified based boundary word with monodromy $h(x)=x+c$, the
boundary cycle containing $x$ is indexed by
$x\bmod\gcd(W,c)$, and its degree is $W/\gcd(W,c)$.
For a boundary reflection $h(x)=-x+c$, the label is the unordered
pair $\{x,c-x\bmod W\}$ and the degree is the size of this pair.
The function \texttt{boundary\_lift\_key} returns this information,
together with the ambient component label.  The chosen based boundary
words specify the paths needed to interpret fibre labels consistently.

The metadata \texttt{cover\_isomorphism\_type\_count} implements
Corollary~\ref{cor:dihedral-exact-types}.  Each family has a
\texttt{type\_id}; the two exceptional records share this identifier
exactly when $m$ is odd.  Their distinct residue records are retained
because they remain distinct components with distinct attachment
labels.  Type identifiers concern isomorphism over the specified base,
which is finer information than abstract genus and boundary count.

Consequently, a polynomially sized list of marked attachment requests
can be resolved in polynomial time without enumerating the unmarked
lifts.  This is a query interface, not a claim that arbitrary
attachments can be forgotten.  If different sheets carry different
external attaching maps, the component-family record alone does not
describe the resulting three-manifold.  That attaching information
must still be retained, either explicitly for finitely many marks or in
a separately justified symbolic representation.

\subsection{Bit complexity and verification}

Let $L$ be the total length of the supplied boundary words, and let
$B$ bound the bit lengths of the integer inputs, $W$, and $r+L+b+1$.
The algorithm performs $O(r+L+b)$ affine operations and gcd/extended-gcd
computations on integers with $O(B)$ bits.  A conservative elementary
bit bound is $O((r+L+b)B^3)$; the purpose of this bound is polynomial
dependence on the binary sheet count, not a sharp arithmetic exponent.
The output has $O((r+L+b)B)$ bits.  For the canonical schema,
$L=O(r+b)$.  The component and boundary queries each take polynomial
time in $B$.

In particular, neither the number of components nor the number of
boundary circles is silently charged as a unary input quantity.  A
reflection cover of a M\"obius band with $W=2^{10000}$ has
$W/2+1$ components, but the implementation returns three records:
$(W-2)/2$ annuli and two M\"obius bands.  A trivial $W$-sheeted cover
of a disk uses one record for $W$ disks.  No loop in the compressed
classifier iterates over $W$ or over its divisors.

The delivered tests compare the arithmetic output against explicit
lifted-spine graph search, parity propagation for orientability, and
direct permutation-cycle enumeration for every boundary circle.  The
independent oracle does not call the classifier's gcd routine, word
evaluator, or surface-schema generator.  There are 13,846 such topology
comparisons: all rank-one assignments through 40 sheets, all rank-two
assignments through 9 sheets, all rank-three assignments through 4
sheets, and 2,000 seeded larger instances, together with targeted
degenerate examples.  The tests also check every component and boundary
query on these instances, 300 simultaneous conjugations of the sheet
labels, malformed inputs, and the $2^{10000}$-sheet examples.
The explicit generic cover-isomorphisms are additionally checked for
equivariance on 156 dihedral presentations.
For the exceptional-isomorphism criterion, 60 small presentations are
checked by enumerating all possible bijections between the two
exceptional fibres and testing the original generating permutations
directly.  These checks do not assume that an equivariant bijection
has to be a translation.  Four further cases test the type counts at
the degenerate moduli $W=1,2$.

\subsection{The precise connection to compressed hierarchies}

Lackenby's compressed-cutting result identifies parallelity-bundle base
topology, bundle type, and attaching information in polynomial time in
the logarithm of surface weight \cite[Theorem~14.4]{Lackenby2026}.
The general orbit-counting work is supplied by
Agol--Hass--Thurston \cite{AHT2006}.  The theorem above supplies a
particularly explicit specialization when a candidate base surface is
already given as a dihedral cover of a smaller bordered surface.
It can answer the topology and marked-query parts of that representation
using binary arithmetic.

This is an actual special case of interval-pairing orbit counting:
on $\{0,\ldots,W-1\}$, translation by $a$ modulo $W$ splits into at
most two ordinary translation pairings, with a break at $W-a$;
reflection $x\mapsto a-x\pmod W$ splits into at most two reversal
pairings, with a break at $a+1$.  The extra hypothesis is that these
pairings come from globally compatible cyclic affine maps.  General
interval pairings do not have that restriction.

There are three remaining obligations before using it inside an unknot
recognizer.  One must extract and verify the covering presentation from
the actual normal-surface or handle data; one must preserve the full
attaching information; and one must prove a global bound on the number
and size of hierarchy states visited.  General normal-surface interval
pairings need not have affine cyclic holonomy, and no theorem here says
that all parallelity bundles possess this additional structure.
Accordingly, this is a proved compressed geometric subroutine and an
explicit prospective fast path.  It is not a proof of a
quasi-polynomial running time for arbitrary unknot diagrams.

